% !TeX root = part01.tex

\documentclass[../final.tex]{subfiles}

\begin{document}

% ============================================================
% Энергетические уровни атомов с одним электроном
% Водородоподобные ионы
% ============================================================

\chapter{Энергетические уровни атомов с одним электроном}

\rsubtitle{Водородоподобные ионы}

Для электрона, движущегося в кулоновском поле ядра с зарядом $Ze$,
стационарное уравнение Шрёдингера имеет вид
%
\begin{equation*}
    \Delta \psi(\mathbf r)
    +
    \frac{2m}{\hbar^2}
    \left(
        W + \frac{Ze^2}{r}
    \right)
    \psi(\mathbf r)
    = 0.
\end{equation*}

Здесь $m$ --- масса электрона, $M$ --- масса ядра.

С учётом конечной массы ядра энергия стационарного состояния равна
%
\begin{equation*}
    W_n
    =
    -\frac{mZ^2e^4}
    {2\hbar^2
    \left(1+\dfrac{m}{M}\right)}
    \frac{1}{n^2}.
\end{equation*}

Так как обычно
%
\[
    M \gg m,
\]
%
можно использовать приближение
%
\begin{equation*}
    W_n
    \approx
    -\frac{mZ^2e^4}{2\hbar^2}
    \frac{1}{n^2}.
\end{equation*}

Введём \rtxt{3}{энергию Ридберга}
%
\begin{equation*}
    Ry
    =
    \frac{me^4}{2\hbar^2}
    \approx
    13.6\,\text{эВ}.
\end{equation*}

Тогда уровни энергии водородоподобного атома удобно записать как

\[
    E_n
    =
    -Ry\,\frac{Z^2}{n^2}.
\]

Таким образом, в нерелятивистском кулоновском приближении энергия
зависит только от главного квантового числа $n$.


% ------------------------------------------------------------
\section{Электронвольт}
% ------------------------------------------------------------

Электронвольт --- единица энергии, равная работе по переносу заряда
одного электрона через разность потенциалов $1\,\text{В}$:
%
\begin{equation*}
    \Delta E
    =
    e\,\Delta\varphi.
\end{equation*}

При $\Delta\varphi=1\,\text{В}$
%
\begin{equation*}
    1\,\text{эВ}
    =
    1.602\cdot10^{-19}\,\text{Кл}\cdot1\,\text{В}
    =
    1.602\cdot10^{-19}\,\text{Дж}.
\end{equation*}

% ------------------------------------------------------------
\section{Волновые функции водородоподобного атома}
% ------------------------------------------------------------

В центрально-симметричном кулоновском поле волновая функция
разделяется на радиальную и угловую части:
%
\begin{equation*}
    \psi_{nlm}(r,\theta,\varphi)
    =
    R_{nl}(r)\,
    Y_{lm}(\theta,\varphi).
\end{equation*}

Сферическая функция также может быть представлена в виде
%
\begin{equation*}
    Y_{lm}(\theta,\varphi)
    =
    \Theta_{lm}(\theta)\,
    \Phi_m(\varphi),
\end{equation*}
%
где зависимость от азимутального угла имеет вид
%
\begin{equation*}
    \Phi_m(\varphi)
    =
    \frac{1}{\sqrt{2\pi}}
    e^{im\varphi}.
\end{equation*}


% ============================================================
\section{Частоты переходов}
% ============================================================

При переходе между двумя стационарными состояниями энергия фотона
равна разности энергий уровней:
%
\begin{equation*}
    \hbar\omega_{i\rightarrow j}
    =
    E_i-E_j.
\end{equation*}

Следовательно,
%
\begin{equation*}
    \omega_{i\rightarrow j}
    =
    \frac{E_i-E_j}{\hbar}.
\end{equation*}

Для водородоподобного атома
%
\[
    \omega_{i\rightarrow j}
    =
    \frac{Ry\,Z^2}{\hbar}
    \left(
        \frac{1}{n_j^2}
        -
        \frac{1}{n_i^2}
    \right).
\]

Для атома водорода $Z=1$, поэтому
%
\begin{equation*}
    \omega_{i\rightarrow j}
    =
    \frac{Ry}{\hbar}
    \left(
        \frac{1}{n_j^2}
        -
        \frac{1}{n_i^2}
    \right).
\end{equation*}

\textit{Замечание.}
В исходном конспекте в этой формуле множитель $Z^2$ не указан,
то есть записан непосредственно случай атома водорода.


% ------------------------------------------------------------
\subsection{Формула Бальмера}
% ------------------------------------------------------------

Для серии Бальмера переходы происходят на уровень $n=2$.
Длины волн спектральных линий могут быть записаны в виде
%
\begin{equation*}
    \lambda
    =
    3645.6\,\text{\AA}\,
    \frac{n^2}{n^2-2^2},
    \qquad
    n=3,4,5,\ldots
\end{equation*}

или
%
\begin{equation*}
    \lambda
    =
    3645.6\,\text{\AA}\,
    \frac{n^2}{n^2-4}.
\end{equation*}


% ============================================================
\section{Релятивистская зависимость энергии от импульса}
% ============================================================

Полная релятивистская энергия свободной частицы удовлетворяет
соотношению
%
\begin{equation*}
    E^2
    =
    p^2c^2+m^2c^4.
\end{equation*}

Поэтому её кинетическая энергия равна
%
\begin{equation*}
    T
    =
    \sqrt{p^2c^2+m^2c^4}
    -
    mc^2
    =
    c\sqrt{p^2+m^2c^2}-mc^2.
\end{equation*}

Нерелятивистская кинетическая энергия
%
\begin{equation*}
    T_0
    =
    \frac{p^2}{2m}.
\end{equation*}

Отличие релятивистской энергии от нерелятивистской будем рассматривать
как малое \rtxt{3}{возмущение}:
%
\begin{equation*}
    V_{\mathrm{rel}}
    =
    T-T_0.
\end{equation*}

При
%
\[
    \frac{p}{mc}\ll1
\]
%
разложим корень:
%
\begin{align*}
    T
    &=
    mc^2
    \left[
        \sqrt{
            1+\frac{p^2}{m^2c^2}
        }
        -1
    \right]
    \\[4pt]
    &=
    mc^2
    \left[
        1
        +
        \frac{1}{2}\frac{p^2}{m^2c^2}
        -
        \frac{1}{8}
        \frac{p^4}{m^4c^4}
        +\ldots
        -1
    \right]
    \\[4pt]
    &=
    \frac{p^2}{2m}
    -
    \frac{p^4}{8m^3c^2}
    +\ldots
\end{align*}

Следовательно,

\[
    V_{\mathrm{rel}}
    =
    -\frac{p^4}{8m^3c^2}.
\]

Это первая релятивистская поправка к кинетической части
гамильтониана.


% ============================================================
\section{Спин-орбитальная связь}
% ============================================================

Для полного момента
%
\begin{equation*}
    \mathbf J
    =
    \mathbf L+\mathbf S
\end{equation*}
%
имеем
%
\begin{align*}
    \mathbf J^2
    &=
    \mathbf L^2
    +
    \mathbf S^2
    +
    2\mathbf L\cdot\mathbf S.
\end{align*}

Отсюда
%
\begin{equation*}
    \mathbf L\cdot\mathbf S
    =
    \frac{1}{2}
    \left(
        \mathbf J^2
        -
        \mathbf L^2
        -
        \mathbf S^2
    \right).
\end{equation*}

На состоянии с определёнными $j,l,s$:
%
\begin{equation*}
    \mathbf L\cdot\mathbf S
    =
    \frac{\hbar^2}{2}
    \left[
        j(j+1)
        -
        l(l+1)
        -
        s(s+1)
    \right].
\end{equation*}

Если угловые моменты выражаются в единицах $\hbar$, то ту же формулу
записывают без явного множителя $\hbar^2$:
%
\begin{equation*}
    \mathbf l\cdot\mathbf s
    =
    \frac{1}{2}
    \left[
        j(j+1)
        -
        l(l+1)
        -
        s(s+1)
    \right].
\end{equation*}


% ------------------------------------------------------------
\subsection{Среднее значение \texorpdfstring{$1/r^3$}{1/r3}}
% ------------------------------------------------------------

Для водородоподобных состояний при $l\neq0$ используется выражение
%
\begin{equation*}
    \left\langle
        \frac{1}{r^3}
    \right\rangle
    =
    \frac{Z^3}{a^3}
    \frac{1}{
        n^3
        l
        \left(l+\dfrac12\right)
        (l+1)
    },
\end{equation*}
%
где $a$ --- боровский радиус.


% ============================================================
\section{Дарвиновская поправка}
% ============================================================

Дарвиновский член имеет вид
%
\[
    V_{\mathrm D}
    =
    \frac{\pi Ze^2\hbar^2}{2m^2c^2}
    \delta(\mathbf r).
\]

Поскольку здесь присутствует $\delta(\mathbf r)$, вклад определяется
значением волновой функции непосредственно в начале координат.

Для $s$-состояний водородоподобного атома
%
\begin{equation*}
    \psi_{n00}(0)
    =
    \frac{1}{\sqrt{\pi}}
    \left(
        \frac{Z}{na}
    \right)^{3/2}.
\end{equation*}

Для состояний с $l\neq0$
%
\begin{equation*}
    \psi_{nlm}(0)=0,
\end{equation*}
%
поэтому дарвиновская поправка существенна только для
\rtxt{3}{$s$-состояний}.


% ============================================================
\section{Изотопический сдвиг}
% ============================================================

При учёте движения ядра вместо массы электрона $m$ в задаче возникает
\rtxt{3}{приведённая масса}
%
\begin{equation*}
    \mu
    =
    \frac{mM}{m+M}.
\end{equation*}

Преобразуем это выражение:
%
\begin{align*}
    \mu
    &=
    m\frac{M}{m+M}
    \\[3pt]
    &=
    m\frac{m+M-m}{m+M}
    \\[3pt]
    &=
    m
    \left(
        1-\frac{m}{m+M}
    \right).
\end{align*}

При $M\gg m$
%
\begin{equation*}
    \frac{m}{m+M}
    \approx
    \frac{m}{M},
\end{equation*}
%
поэтому

\[
    \mu
    \approx
    m
    \left(
        1-\frac{m}{M}
    \right).
\]

Энергетические уровни при этом равны
%
\begin{equation*}
    E_n
    =
    -Z^2
    \frac{\mu e^4}{2\hbar^2}
    \frac{1}{n^2},
\end{equation*}
%
или в первом приближении по отношению $m/M$
%
\begin{equation*}
    E_n
    \approx
    -Z^2
    \frac{me^4}{2\hbar^2}
    \frac{1}{n^2}
    \left(
        1-\frac{m}{M}
    \right).
\end{equation*}

Таким образом, уровни энергии слабо зависят от массы ядра $M$.
У разных изотопов одного элемента массы ядер различаются, поэтому
слегка различаются и положения спектральных линий. Это явление
называется \rtxt{3}{изотопическим сдвигом}.


% ------------------------------------------------------------
\subsection{Постоянная Ридберга}
% ------------------------------------------------------------

Для бесконечно тяжёлого ядра
%
\begin{equation*}
    R_{\infty}
    \approx
    109737.31\,\text{см}^{-1}.
\end{equation*}

В спектроскопии часто используется волновое число
%
\begin{equation*}
    \widetilde{\nu}
    =
    \frac{1}{\lambda}.
\end{equation*}

Если длина волны $\lambda$ выражена в сантиметрах, то
%
\begin{equation*}
    [\widetilde{\nu}]
    =
    \text{см}^{-1}.
\end{equation*}

\end{document}